% Differential-evolution baseline: exact description of the implementation
% (src/hsearch/search/strategies.py, class DEZelinka; harness in src/hsearch/search/harness.py;
% configuration in experiments/e50_batch_runner.py and e20_search_benchmark.py).
% One paragraph, no table; intended for the appendix.

\textbf{Differential-evolution baseline.} The baseline is differential evolution over initial conditions with the cost function of Ref.~\onlinecite{Zelinka2016}, in the variant DE/rand/1/bin used there. The population has $N_P=25$ members (half of the initial sample $n_0=50$ of the other strategies, but not fewer than 10), the differential weight is $F=0.9$, and the crossover probability is $C_R=0.3$, the values of Ref.~\onlinecite{Zelinka2016}. The initial population is drawn uniformly from the region $B$ and consists of the first 25 points of the uniform stream of the run, so that it coincides with the first 25 points of uniform sampling (common random numbers). In every generation, for each member $x_i$, three other members $x_{r_1}$, $x_{r_2}$, $x_{r_3}$, distinct from one another and from $x_i$, are drawn without replacement; the mutant is $v=x_{r_3}+F\,(x_{r_1}-x_{r_2})$; the trial $u_i$ takes each coordinate from $v$ with probability $C_R$ and from $x_i$ otherwise, with one randomly chosen coordinate always taken from $v$; and $u_i$ is clipped to $B$ coordinate by coordinate. The 25 trials of a generation form one batch and are integrated with the full protocol ($T_{\rm tr}+T_{\rm obs}$ of Table~\ref{tab:integration}); $u_i$ replaces $x_i$ if its cost is not larger than that of $x_i$. The cost is computed from the trajectory fingerprint alone: with $r_{\max}$ the maximum of the Euclidean norm of the state over the observation window $T_{\rm obs}$ (after the transient $T_{\rm tr}$), the cost is 0 if $d_{\rm lo}\le r_{\max}\le d_{\rm hi}$ and $r_{\max}$ otherwise, and it is $10^6$ for a trajectory that escapes beyond $R_{\rm esc}$ or produces a non-finite state at any time. This is the function of Ref.~\onlinecite{Zelinka2016}, where the maximum is taken over the last 50 time units of a run of length 200 and the interval $[d_{\rm lo},d_{\rm hi}]=[3,11]$ was chosen empirically for the hidden attractor of the Chua circuit; the pair $(3,11)$ remains the default of the implementation but is not used in the benchmark. Instead, the interval is set from the size of the target, which is the sense in which the baseline is ``told the size'': if $r^\ast$ denotes the value of $r_{\max}$ in the fingerprint of the target reference (the mean over its 64 calibration copies; for a problem with several target regimes, the largest of their values), then $[d_{\rm lo},d_{\rm hi}]=[0.5\,r^\ast,\,1.5\,r^\ast]$. For Chua~I, $r^\ast=10.9$ and the interval is $[5.5,16.4]$. The cost does not use the labels of the identification protocol, so the search does not distinguish the target from any other bounded regime whose size falls into the interval. Budget and detection are the same as for every other strategy: each integrated trial counts as one evaluation, the initial population included; a run has the per-problem budget of Sec.~\ref{sec:measure} (5--15 times $N^0_{0.9}$ on levels A and B, from $10^3$ integrations on Chua~I to $4\times10^4$ on the Jafari chaos); and the first hit is the ordinal number, in the order of evaluation, of the first integrated point whose fingerprint the identification protocol matches to the target reference, whether or not that point is accepted into the population. The run stops at the end of the generation in which the hit occurs, and $N_{0.9}$ is the 0.9 quantile of the first-hit number over the $R$ runs.
